\documentclass[11pt]{scrartcl}

\usepackage[top=1.5cm]{geometry}
\usepackage{url}
\usepackage{float}
\usepackage{listings}
\usepackage{xcolor}
\usepackage{graphicx}
\usepackage{booktabs}

\setlength{\parindent}{0em}
\setlength{\parskip}{0.5em}

\renewcommand{\thesubsection}{\arabic{subsection}}

\lstdefinestyle{dmrsql}{
  language=SQL,
  basicstyle=\small\ttfamily,
  keywordstyle=\color{magenta!75!black},
  stringstyle=\color{green!50!black},
  showspaces=false,
  showstringspaces=false,
  commentstyle=\color{gray}}

\lstdefinestyle{dmrPython}{
  language=Python,
  basicstyle=\small\ttfamily,
  keywordstyle=\color{magenta!75!black},
  stringstyle=\color{green!50!black},
  showspaces=false,
  showstringspaces=false,
  commentstyle=\color{gray}}

\title{
  \textbf{\large Projektaufgabe 3 } \\
  Phase 1 -- Das EDGE Modell (4 P) \\
  {\large Datenmanagement jenseits von Relationen}
}

\author{
  TODO: Gruppen Nummer \\
  \large TODO: Lastname1 Firstname1, StudentID1 \\
  \large TODO: Lastname2 Firstname2, StudentID2
}

\begin{document}

\maketitle\thispagestyle{empty}

Dieses Reporting Template dient der Vorbereitung der Abgabe von Phase 1.

\subsection*{EDGE Modell Import des Toy Beispiels (1 Punkte)}

Beschreiben Sie kurz, welches Verfahren (inkl. Programmiersprache) Sie fuer den Import verwenden. Gab es Probleme beim Implementieren des Parsers? (0.5 P).

Der Import wurde in Python 3 umgesetzt. Das Toy-Beispiel wird mit \texttt{xml.etree.ElementTree} geparst und danach in eine eigene Baumstruktur \texttt{EdgeNode} ueberfuehrt. In diesem Schritt wird die im Aufgabenblatt geforderte Transformation ausgefuehrt: Die Wurzel bleibt \texttt{bib}; darunter werden die Publikationen zuerst nach Venue und danach nach Jahr gruppiert. Dadurch entsteht die Struktur
\[
\texttt{bib} \rightarrow \texttt{venue} \rightarrow \texttt{year} \rightarrow \texttt{publication} \rightarrow \texttt{fields}.
\]

Die Venue wird aus dem \texttt{key}-Attribut der DBLP-Publikation abgeleitet. Die Prefixe \texttt{journals/pvldb/} und \texttt{conf/vldb/} werden als \texttt{vldb} interpretiert, \texttt{journals/pacmmod/} und \texttt{conf/sigmod/} als \texttt{sigmod}, und \texttt{conf/icde/} als \texttt{icde}. Das Jahr wird aus dem Kindknoten \texttt{year} gelesen. Die Attribute \texttt{mdate} und \texttt{orcid} werden gemaess Aufgabenstellung ignoriert.

Ein praktisches Problem waren HTML/XML-Entities im DBLP-Ausschnitt, z.B. \texttt{\&uuml;} und \texttt{\&auml;}. Vor dem Parsen werden deshalb bekannte benannte Entities ueber die Python-Standardbibliothek \texttt{html.entities} auf Unicode-Zeichen abgebildet. Die XML-Standardentities wie \texttt{\&amp;} bleiben dabei erhalten.

Zeigen Sie (z.B. mittels Screenshot des Konsolenoutputs), dass Ihr Parser alle Elemente (und deren Inhalt) des Toy-Beispiels korrekt erkennt. (0.5 P).

Der Parser kann mit folgendem Befehl ausgefuehrt werden:

\begin{lstlisting}[style=dmrPython]
python .\projektaufgabe_3\phase1_demo.py
\end{lstlisting}

Der Konsolenoutput zeigt die transformierte Baumstruktur. Ein Ausschnitt ist:

\begin{lstlisting}
0: s_id=bib, type=bib, content=NULL
  1: s_id=vldb, type=venue, content=NULL
    2: s_id=vldb_2023, type=year, content=NULL
      3: s_id=SchmittKAMM23, type=article, content=NULL
        4: s_id=NULL, type=author, content=Daniel Ulrich Schmitt
        5: s_id=NULL, type=author, content=Daniel Kocher
        6: s_id=NULL, type=author, content=Nikolaus Augsten
        9: s_id=NULL, type=title,
           content=A Two-Level Signature Scheme for Stable Set Similarity Joins.
      18: s_id=SchalerHS23, type=article, content=NULL
        19: s_id=NULL, type=author, content=Christine Schaeler
  31: s_id=sigmod, type=venue, content=NULL
    32: s_id=sigmod_2022, type=year, content=NULL
      33: s_id=HutterAK0L22, type=inproceedings, content=NULL
    46: s_id=sigmod_2023, type=year, content=NULL
      47: s_id=ThielKAHMS23, type=article, content=NULL
\end{lstlisting}

Der vollstaendige Output enthaelt alle vier Publikationen des Toy-Beispiels mit ihren Feldknoten, z.B. \texttt{author}, \texttt{title}, \texttt{pages}, \texttt{year}, \texttt{journal}, \texttt{ee} und \texttt{url}.

Geben Sie die Definition Ihrer Tabellen im Edge Modell an. Zeigen Sie (einen Ausschnitt) des Toy Beispiels nach dem Import als Edge Model in die DB (1 P).

Wir interpretieren das EDGE Modell als zwei Relationen. Die Relation \texttt{node} speichert alle Knoten des transformierten XML-Baums. \texttt{id} ist der technische Knotenschluessel, \texttt{s\_id} ist ein semantischer Bezeichner fuer Wurzel, Venue-, Year- und Publikationsknoten, \texttt{type} enthaelt den Knotentyp, und \texttt{content} enthaelt den Textinhalt von Blattknoten. Die Relation \texttt{edge} speichert gerichtete Eltern-Kind-Kanten.

\begin{lstlisting}[style=dmrsql]
CREATE TABLE IF NOT EXISTS node (
    id INT PRIMARY KEY,
    s_id TEXT,
    type TEXT NOT NULL,
    content TEXT
);

CREATE TABLE IF NOT EXISTS edge (
    from_id INT NOT NULL REFERENCES node(id),
    to_id INT NOT NULL REFERENCES node(id),
    PRIMARY KEY (from_id, to_id)
);

CREATE INDEX IF NOT EXISTS idx_node_s_id ON node(s_id);
CREATE INDEX IF NOT EXISTS idx_node_type ON node(type);
CREATE INDEX IF NOT EXISTS idx_node_content ON node(content);
CREATE INDEX IF NOT EXISTS idx_edge_from ON edge(from_id);
CREATE INDEX IF NOT EXISTS idx_edge_to ON edge(to_id);
\end{lstlisting}

Nach dem Import des Toy-Beispiels enthaelt die Datenbank 62 Tupel in \texttt{node} und 61 Tupel in \texttt{edge}. Ein Ausschnitt der Relation \texttt{node} ist:

\begin{table}[H]
  \centering
  \begin{tabular}{rlll}
    \toprule
    id & s\_id & type & content \\
    \midrule
    0 & bib & bib & NULL \\
    1 & vldb & venue & NULL \\
    2 & vldb\_2023 & year & NULL \\
    3 & SchmittKAMM23 & article & NULL \\
    4 & NULL & author & Daniel Ulrich Schmitt \\
    5 & NULL & author & Daniel Kocher \\
    6 & NULL & author & Nikolaus Augsten \\
    9 & NULL & title & A Two-Level Signature Scheme \ldots \\
    18 & SchalerHS23 & article & NULL \\
    31 & sigmod & venue & NULL \\
    32 & sigmod\_2022 & year & NULL \\
    33 & HutterAK0L22 & inproceedings & NULL \\
    46 & sigmod\_2023 & year & NULL \\
    47 & ThielKAHMS23 & article & NULL \\
    \bottomrule
  \end{tabular}
  \caption{Ausschnitt der Relation \texttt{node} nach Import des Toy-Beispiels}
\end{table}

Ein Ausschnitt der Relation \texttt{edge} ist:

\begin{table}[H]
  \centering
  \begin{tabular}{rr}
    \toprule
    from\_id & to\_id \\
    \midrule
    0 & 1 \\
    1 & 2 \\
    2 & 3 \\
    3 & 4 \\
    3 & 5 \\
    3 & 6 \\
    2 & 18 \\
    0 & 31 \\
    31 & 32 \\
    32 & 33 \\
    31 & 46 \\
    46 & 47 \\
    \bottomrule
  \end{tabular}
  \caption{Ausschnitt der Relation \texttt{edge} nach Import des Toy-Beispiels}
\end{table}

\subsection*{Berechnung der XPath Achsen im EDGE Modell (2 Punkte)}

Die vier XPath-Achsen werden auf dem EDGE Modell berechnet. Die rekursiven Achsen \texttt{ancestor} und \texttt{descendant} werden in SQL mit \texttt{WITH RECURSIVE} und Joins ueber die Relation \texttt{edge} berechnet. Die Achsen \texttt{following-sibling} und \texttt{preceding-sibling} werden berechnet, indem zuerst der gemeinsame Elternknoten ueber \texttt{edge} bestimmt wird und danach die Geschwister mit groesserer bzw. kleinerer ID ausgegeben werden.

Geben Sie fuer jede der geforderten Achsen, die Knotenmenge des Ergebnisses, sowie die Anzahl der Ergebnisknoten aus.

\begin{table}[h]
	\centering
		\begin{center}
			\begin{tabular}{ l | c c }
				\toprule
				Achse & Ergebnisknoten ID's & Ergebnisgroesse\\
				\midrule
				ancestor & 0, 1, 2, 3 & 4 \\
				descendants & 3--30 & 28 \\
				following SchmittKAMM23 & 18 & 1 \\
				preceeding SchmittKAMM23 & $\emptyset$ & 0 \\
				following SchalerHS23 & $\emptyset$ & 0 \\
				preceeding SchalerHS23 & 3 & 1 \\
				\bottomrule
			\end{tabular}
			\end{center}
	\caption{Ergebnisse der XPath-Berechnug}
	\label{tab:ErgebnisseDerXPathBerechnug}
\end{table}

Die Ergebnisse stammen aus folgendem Setup-Lauf:

\begin{lstlisting}
[DB] Saved EDGE model:
node rows: 62
edge rows: 61

[AXES] Toy correctness checks:
ancestor("Daniel Ulrich Schmitt")
  id=0, s_id=bib, type=bib, content=None
  id=1, s_id=vldb, type=venue, content=None
  id=2, s_id=vldb_2023, type=year, content=None
  id=3, s_id=SchmittKAMM23, type=article, content=None

descendant(vldb_2023)
  id=3, s_id=SchmittKAMM23, type=article, content=None
  ...
  id=30, s_id=None, type=url,
     content=db/journals/pvldb/pvldb16.html#SchalerHS23

following-sibling(SchmittKAMM23)
  id=18, s_id=SchalerHS23, type=article, content=None
preceding-sibling(SchmittKAMM23)
  <empty>
following-sibling(SchalerHS23)
  <empty>
preceding-sibling(SchalerHS23)
  id=3, s_id=SchmittKAMM23, type=article, content=None
\end{lstlisting}

Sofern Sie Annahmen zum Aufbau der Daten getroffen haben, geben Sie diese nachfolgend an:
\begin{itemize}
	\item Die Venue-Zuordnung erfolgt ueber das \texttt{key}-Attribut der DBLP-Publikation.
	\item Jede importierte Publikation besitzt genau einen \texttt{year}-Kindknoten, der fuer die Gruppierung verwendet wird.
	\item Die IDs werden nach der Transformation in Preorder-Reihenfolge des EDGE-Baums vergeben. Dadurch entsprechen kleinere IDs bei Geschwistern der vorherigen Dokument-/Baumreihenfolge.
	\item XML-Attribute werden nicht als eigene EDGE-Knoten gespeichert. Die laut Aufgabenstellung ignorierbaren Attribute \texttt{mdate} und \texttt{orcid} werden daher verworfen.
	\item Die Achsenberechnung findet in der Datenbank statt; Python dient nur als Controller fuer die SQL-Anfragen.
\end{itemize}

\subsection*{Zeitmamagement}

Benoetigte Zeit pro Person (nur Phase 1): \textbf{TODO}


\end{document}
